\chapter{สรุปผล อภิปรายผล และข้อเสนอแนะ}

\section{สรุปผลการวิจัย}
การวิจัยนี้มีวัตถุประสงค์เพื่อ (1) ประเมินผลกระทบของโครงการ Smart Farmer ต่อการพัฒนาทักษะของเกษตรกร (2) วิเคราะห์ความสัมพันธ์ระหว่างปัจจัยที่เกี่ยวข้องกับการเปลี่ยนแปลงทักษะ และ (3) เสนอกรอบแนวทางการพัฒนาทักษะเกษตรกรโดยใช้เทคโนโลยี Machine Learning

ทั้งนี้ การสรุปผลในบทนี้จะจัดเรียงให้ตอบวัตถุประสงค์ดังกล่าวโดยเชื่อมโยงกับผลการวิเคราะห์เชิงประจักษ์ในบทที่ 4 อย่างเป็นระบบ

\begin{enumerate}

\item \textbf{ผลกระทบของโครงการ Smart Farmer ต่อการพัฒนาทักษะของเกษตรกร}

จากการประเมินผลกระทบด้วยวิธี Propensity Score Matching ร่วมกับ Difference-in-Differences (PSM-DID) ในบทที่ 4 พบว่า โดยภาพรวมการเข้าร่วมโครงการไม่ทำให้คะแนนทักษะของเกษตรกรเพิ่มขึ้นอย่างมีนัยสำคัญทางสถิติเมื่อพิจารณาในระดับค่าเฉลี่ย

ผลดังกล่าวแสดงให้เห็นว่า ผลลัพธ์ของโครงการอาจไม่เด่นชัดในภาพรวมภายใต้บริบทที่ผู้เข้าร่วมมีความแตกต่างกันสูง รวมถึงมีข้อจำกัดด้านการนำความรู้ไปประยุกต์ใช้ในภาคปฏิบัติ และสะท้อนว่า รูปแบบการดำเนินโครงการในปัจจุบันอาจยังไม่สามารถสร้างการเปลี่ยนแปลงเชิงพฤติกรรมหรือทักษะที่แตกต่างจากแนวโน้มทั่วไปได้อย่างมีนัยสำคัญ

อย่างไรก็ดี ผลการประเมินเชิงสาเหตุด้วย PSM-DID ยังมีความสำคัญในฐานะหลักฐานที่ชี้ว่า เมื่อควบคุมความแตกต่างเชิงลักษณะพื้นฐานระหว่างกลุ่มผู้เข้าร่วมและกลุ่มเปรียบเทียบแล้ว ผลสัมฤทธิ์ด้านทักษะโดยเฉลี่ยยังไม่ปรากฏเด่นชัด

\item \textbf{ความสัมพันธ์ระหว่างปัจจัยที่เกี่ยวข้องกับการเปลี่ยนแปลงทักษะ}

การวิเคราะห์เพิ่มเติมด้วยเทคนิค Machine Learning และการอธิบายตัวแบบด้วย SHAP พบว่า ปัจจัยที่สัมพันธ์กับการเปลี่ยนแปลงทักษะมีลักษณะเป็น “ปัจจัยร่วม” ที่เชื่อมโยงทั้งความพร้อมของเกษตรกรและเงื่อนไขแวดล้อมของพื้นที่

โดยปัจจัยสำคัญประกอบด้วย

\begin{itemize}
\item ระดับทักษะเดิม (baseline skills)
\item ทักษะด้านเทคโนโลยีและระดับการศึกษา
\item โครงสร้างพื้นฐานที่เอื้อต่อการผลิต เช่น ระบบชลประทาน
\end{itemize}

ข้อค้นพบนี้ช่วยขยายความหมายของผล PSM-DID โดยชี้ว่า แม้ผลกระทบเฉลี่ยไม่เด่นชัด แต่ผลลัพธ์มีความแตกต่างกันตามลักษณะของเกษตรกรและบริบทพื้นที่ และสะท้อนลักษณะความไม่เป็นเนื้อเดียวกันของผลลัพธ์ (heterogeneity)

\item \textbf{กรอบแนวทางการพัฒนาทักษะเกษตรกรโดยใช้ Machine Learning}

เมื่อพิจารณาผลการอธิบายตัวแบบ (SHAP) ร่วมกับผลการจัดกลุ่ม (clustering) พบว่า เกษตรกรสามารถจำแนกออกเป็นหลายกลุ่มที่มีศักยภาพและข้อจำกัดแตกต่างกันอย่างชัดเจน

จึงนำไปสู่ข้อเสนอเชิงกรอบแนวทางว่า การพัฒนาทักษะไม่ควรใช้แนวทางแบบเดียวกับทุกกลุ่ม (one-size-fits-all) แต่ควรออกแบบการสนับสนุนแบบจำเพาะกลุ่ม (targeted intervention)

โดยใช้ข้อมูลเป็นฐานเพื่อ

\begin{itemize}
\item คัดกรอง/จัดกลุ่มเกษตรกรตามความพร้อมและข้อจำกัด
\item จัดสรรรูปแบบการอบรมและการติดตามผลให้เหมาะสมกับแต่ละกลุ่ม
\item ใช้โมเดล Machine Learning เพื่อคาดการณ์โอกาสเกิดผลลัพธ์และสนับสนุนการตัดสินใจเชิงนโยบาย
\end{itemize}

\end{enumerate}

\section{อภิปรายผลการวิจัย (Discussion)}

\subsection{การตีความผลการประเมินผลกระทบ (DID)}

หัวข้อนี้มุ่งอธิบายความหมายของข้อค้นพบในบทที่ 4 โดยเชื่อมโยงผลการประเมินเชิงสาเหตุด้วย PSM-DID เข้ากับผลการวิเคราะห์เชิงลึกด้วย Machine Learning/SHAP และผลการจัดกลุ่ม (clustering)

จากข้อค้นพบดังกล่าว สามารถอภิปรายประเด็นหลักได้ 2 ประการ ได้แก่ (1) ปัญหา Selection Bias และ (2) ข้อจำกัดของการวัดผลในระดับค่าเฉลี่ย

\textbf{(1) ปัญหา Selection Bias}

ผลการเปรียบเทียบลักษณะพื้นฐานในบทที่ 4 บ่งชี้ว่า เกษตรกรที่เข้าร่วมโครงการมีแนวโน้มเป็นกลุ่มที่มีความพร้อมสูงกว่าอยู่แล้ว แม้วิธี PSM จะช่วยลดความเอนเอียงดังกล่าวได้บางส่วน แต่ปัจจัยที่ไม่สามารถสังเกตได้ยังอาจส่งผลต่อค่าประมาณ

\textbf{(2) ข้อจำกัดของการวัดผลในระดับค่าเฉลี่ย}

เมื่อผลลัพธ์มีความแตกต่างกันระหว่างกลุ่ม การใช้ค่าเฉลี่ยเพียงค่าเดียวอาจทำให้ผลกระทบที่เกิดกับบางกลุ่มถูกกลบ ดังนั้น การใช้ Machine Learning และ SHAP จึงช่วยเปิดเผยความแตกต่างในระดับรายบุคคลได้ชัดเจนขึ้น

\subsection{บทบาทของปัจจัยโครงสร้างและทักษะเดิม}

ผลการวิเคราะห์ชี้ว่า ปัจจัยพื้นฐาน เช่น ทักษะเดิม การศึกษา และทรัพยากร มีบทบาทสำคัญต่อผลลัพธ์ สะท้อนว่า การอบรมเพียงอย่างเดียวอาจไม่เพียงพอ

\subsection{นัยเชิงระบบ (Systemic Insight)}

เมื่อพิจารณาผลทั้งหมดร่วมกันสามารถสรุปเชิงระบบได้ว่า ข้อจำกัดของผลลัพธ์โครงการไม่ได้อยู่ที่ “เนื้อหาการอบรม” เพียงอย่างเดียว แต่เกี่ยวข้องกับการออกแบบมาตรการและกลไกสนับสนุนที่ยังไม่สอดรับกับความแตกต่างของผู้เข้าร่วม

ดังนั้น การพัฒนาโครงการควรมุ่งสร้างรูปแบบการสนับสนุนที่เหมาะกับแต่ละกลุ่ม และเพิ่มเงื่อนไขเอื้อเพื่อให้เกิดการนำความรู้ไปใช้จริง

\section{ข้อเสนอแนะเชิงนโยบาย (Policy Recommendations)}

\subsection{การคัดกรองและจัดกลุ่มเป้าหมายบนฐานข้อมูล}

ควรใช้ตัวแบบ Machine Learning เพื่อคัดกรองและจัดกลุ่มเกษตรกรตามระดับความพร้อม โดยพิจารณาจากข้อมูลพื้นฐานและปัจจัยสำคัญ

\subsection{การออกแบบหลักสูตรแบบจำเพาะกลุ่ม}

\begin{itemize}
\item กลุ่มพร้อมสูง: หลักสูตรขั้นสูงและการใช้เทคโนโลยี
\item กลุ่มกลาง: เสริมทักษะดิจิทัลและการนำไปใช้
\item กลุ่มข้อจำกัด: เริ่มจากพื้นฐานและทรัพยากร
\end{itemize}

\subsection{การพัฒนาโครงสร้างพื้นฐาน}

ควรพัฒนาโครงสร้างพื้นฐาน เช่น ระบบชลประทาน เทคโนโลยี และบริการสนับสนุนควบคู่การอบรม

\subsection{การติดตามผลระยะยาว}

ควรกำหนดระบบติดตามผลแบบ longitudinal เพื่อประเมินผลกระทบที่แท้จริง

\section{ข้อเสนอแนะสำหรับการวิจัยในอนาคต (Future Research)}

\subsection{การศึกษาระยะยาว}
ควรมี longitudinal study เพื่อวัดผลระยะยาว

\subsection{การลด Selection Bias}
ควรใช้ RCT หรือ quasi-experimental design ที่เข้มแข็งขึ้น

\subsection{การพัฒนา Recommender System}
พัฒนาระบบแนะนำเพื่อจับคู่เกษตรกรกับหลักสูตร

\subsection{ข้อมูลเชิงพื้นที่}
ควรใช้ spatial data เพื่อวิเคราะห์ความแตกต่าง

\section{ข้อสรุปเชิงนโยบาย (Final Insight)}

การพัฒนาเกษตรกรควรเปลี่ยนจาก one-size-fits-all ไปสู่ targeted policy

Machine Learning ควรถูกใช้เป็นเครื่องมือสนับสนุนการตัดสินใจ

ประสิทธิผลของโครงการขึ้นอยู่กับความสอดคล้องระหว่างการอบรมและความพร้อมของเกษตรกร

การละเลยปัจจัยดังกล่าวนำไปสู่การที่ผลลัพธ์ไม่ปรากฏในระดับค่าเฉลี่ย แม้จะมีผลในบางกลุ่ม